\section{Introduction}
\label{sec:introduction}

A bank validates a value-at-risk (VaR) model by counting the days on which the loss exceeded the forecast, and a rating system by comparing the
defaults of each grade with the probability of default (PD) the grade was assigned. The tests that carry these comparisons were designed to be read
once, on a window fixed in advance: the proportion-of-failures test of \citet{kupiec1995} on a year of daily forecasts, the three-zone traffic light of the Basel
Committee \citep{bcbs1996,bcbs2019} on the last 250 days, the binomial and chi-square tests of the validation studies of the Committee \citep{bcbs2005,tasche2008}
on the defaults of a cohort. A risk function does not read them once. It reads them every day or every quarter, on a growing or a sliding window, and
acts on the first reading that crosses a threshold. A test that has a size of 5\,\% when read once has a larger size when read repeatedly on the same
data, because each reading is another chance to reject a correct model \citep{armitage1969}. On the generator of known truth used in this study, a correct 99\,\% VaR
model is rejected at some day of 1,000 by the Kupiec test on the last 250 days in \nNullSizeKupTrail\,\% of the series, by the same test on all days so far in
\nNullSizeKupCum\,\%, and the traffic light leaves its green zone in \nNullSizeLight\,\% (\cref{tab:main-results}). The nominal 5\,\% describes one reading,
not the practice.

The remedy that the literature on sequential analysis offers is to change the test, not the practice. Group-sequential boundaries \citep{pocock1977,obrien1979,lan1983}
spend a fixed size over a number of looks fixed in advance. Sequential probability ratio tests \citep{wald1945} and confidence sequences \citep{darling1967,howard2021} remove the
fixed number: they keep their guarantee at every time. The unifying object is the e-process, a non-negative process whose expectation under the null does not
exceed one at any stopping time, and which therefore has a size of at most $a$ when the null is rejected the first time it reaches $1/a$, however often it
is read and whenever the analyst stops \citep{shafer2011,shafer2021,ramdas2023,grunwald2024,vovk2021}. \citet{wang2022} built e-processes for the back-testing of VaR and expected shortfall
with a betting rule that learns the size of the departure from the data.

Two gaps remain for a validation function. The first is the calibration of probabilities of default by rating grade: the tests in use are read at the end of a
cohort's life. E-values for the calibration of forecasts exist \citep{arnold2023}, and so do e-values for the comparison of two probability forecasts \citep{henzi2022}, but we found no e-process for the defaults of each rating grade against its predicted PD, with a multiplicative miscalibration as the alternative and batteries of grades, among the works cited here. The second is the combination of several monitors of one portfolio into one anytime-valid decision, where
the combination rule decides whether the guarantee survives \citep{vovk2021,ramdas2023}. We address both and measure what the guarantee costs.

The proposed method is a set of four monitors, each a mixture of likelihood-ratio e-processes over a grid of alternatives fixed before any outcome is seen: a monitor of the
VaR exceedance rate, a monitor of hits that follow hits, a monitor of the defaults of each rating grade against its predicted PD, and a battery that combines monitors by
averaging, which is valid whatever the dependence between them, or by product, which is valid when they are independent given the past. Each monitor rejects the first
time its process reaches 20, which is an anytime-valid size of 5\,\%. The monitors are the module \texttt{riskvalidation.validation.sequential} of the
open engine \texttt{riskvalidation} release 1.1.0 and serve three cases of the validation suite: the sizes and powers of validation tests on planted defects, market-risk
back-testing across four crises, and small-business credit by vintage.

The contributions are four.
\begin{enumerate}
\item The monitors of the VaR exceedance rate, of hit clustering and of the defaults of each grade, with the combination rules of a battery, specified in full, with their
parameters, their algorithm and their cost (\cref{sec:proposed-method}).
\item A theory of what the monitors guarantee and what they cost: validity at any stopping time, validity of the clustering and PD factors, the failure of the
product combination under dependence, a bound on the growth rate of any e-process and the first-order delay it implies, a bound on the power of any monitor at a fixed horizon, and the
loss of validity of the PD monitor when the default probability of a year is random (\cref{sec:theory}).
\item An experiment in which the monitors are run against ten other rules on the same streams: the Kupiec test on a growing and on a sliding window, the traffic light,
Christoffersen's independence test, a Sidak-corrected Kupiec test, a Kupiec test and a CUSUM chart whose thresholds are calibrated by simulation to the same size, the published
e-backtest, and two e-processes of the textbook kind (\cref{sec:experimental-design,sec:results}).
\item The measured outcome of the hypotheses stated before the experiments. The size holds, and the delay and the power cost hold at a hit rate four times too high and above. The delay and the power-cost hypotheses fail at three times the nominal rate: the delay is \nSeedsThreeMed\ days against the window of 250 (\cref{tab:x3-seeds}) and the power cost is \nSeedsThreeCost\ points against the 10 allowed. The PD monitor loses its size once defaults share a systematic factor, and at an equal size it detects a miscalibration of the PD later than a chi-square rule does, where the engine's Vasicek-corrected binomial test with the right correlation keeps its size
(\cref{tab:hypotheses,tab:pd-correlation,tab:pd-comparators,tab:pd-equal-size}).
\end{enumerate}

The main result is a trade. Read every day over 1,000 days, the rate e-process rejects a correct model in \nNullSizeE\,\% of \nReps\ series (95\,\% interval \nNullSizeELo\ to \nNullSizeEHi\,\%)
where the repeated Kupiec test and the traffic light reject in \nNullSizeKupCum\ to \nNullSizeKupTrail\,\%, and it detects a hit rate four times too high at a median of
\nMedFourE\ days, inside the 250-day window that the traffic light needs before it can speak at all. The price is power at a fixed horizon: by day 250 it detects a hit rate
three times too high in \nByThreeE\,\% of the series against \nShotThree\,\% for one Kupiec reading. A one-sided monitor of the direction a supervisor penalises cuts that cost to \nCostThreeEUp\ points in the same run
(\cref{sec:ablation}); that variant was chosen after the two-sided results were known and needs a prospective confirmation (\cref{sec:discussion-negative}).

\Cref{sec:related-work} places the work among the published methods, \cref{sec:problem} states the problem, \cref{sec:proposed-method} the monitors and \cref{sec:theory} their theory.
\Cref{sec:experimental-design} sets the hypotheses and the protocol, \cref{sec:results,sec:ablation} give the measurements, and \cref{sec:discussion,sec:limitations}
read them and bound them. \Cref{sec:reproducibility} lists the commands that regenerate every number.
